\documentclass{article}
% Required packages
\usepackage{amssymb}
\usepackage{amsmath}
\usepackage{graphicx}
\usepackage{geometry}
\usepackage{tikz}
\usepackage{array}
\usepackage{booktabs}
\usepackage{enumitem}
\usepackage{listings}
\usepackage{xcolor}
\usepackage{fancyhdr}
\usepackage{float}
\usepackage{subcaption}
\usepackage{comment}
\usepackage{algorithm}
\usepackage{algpseudocode}
\usepackage{caption}
\usepackage{siunitx}
\usepackage{hyperref}
\usetikzlibrary{arrows.meta, positioning}

% Set page geometry
\geometry{a4paper, margin=1in}

% Configure listings for Python
\lstset{
  language=Python,
  basicstyle=\ttfamily\footnotesize,
  numbers=left,
  numberstyle=\tiny\color{gray},
  frame=single,
  breaklines=true,
  breakatwhitespace=true,
  captionpos=b,
  tabsize=4,
  showspaces=false,
  showstringspaces=false,
  showtabs=false,
  commentstyle=\color{gray}\textit,
  keywordstyle=\color{blue}\bfseries,
  stringstyle=\color{red}
}

% Header setup
\pagestyle{fancy}
\chead{DSC 257: Unsupervised Learning (Fall 2025)}
\lhead{Homework 10}
\rhead{Randall Rogers}

\begin{document}

%------------------
% Solution for 1(a)
%------------------
\subsection*{Solution 1}
\noindent\rule{\textwidth}{0.4pt}\\
\subsubsection*{Solution 1 (a)}
\subsubsection*{Python Code}
\begin{lstlisting}
# Train with latent_dim = 4
latent_dim = 4  # CHANGED FROM 8 TO 4
model = VAE(x_dim=784, z_dim=latent_dim).to(device)
optimizer = optim.Adam(model.parameters())

num_epochs = 10 # 10-20 epochs is usually sufficient for a quick result
for epoch in range(num_epochs):
    epoch_loss = 0
    for batch_idx, (data, _) in enumerate(train_loader):
        data = data.to(device)
        optimizer.zero_grad()
        recon_data, mu, logvar = model(data)
        loss = loss_function(recon_data, data, mu, logvar)
        loss.backward()
        optimizer.step()
        epoch_loss += loss.item()
    print(f"Epoch {epoch+1} Loss: {epoch_loss / len(train_dataset):.4f}")
# Save the model
torch.save(model.state_dict(), f'vae_mnist_z_{latent_dim}.pth')
# Generate 100 Samples
with torch.no_grad():
    # Sample 100 latent vectors from N(0, I)
    random_sample = torch.randn(100, latent_dim).to(device)
    # Decode them
    generated_image = model.decoder(random_sample)
    # Save as a 10x10 grid
    save_image(generated_image.view(100, 1, 28, 28), 
               f'samples_z_{latent_dim}.png', nrow=10)
    print(f"Saved generated samples to samples_z_{latent_dim}.png")
\end{lstlisting}

\subsubsection*{Generated Image}
\begin{center}
          \includegraphics[scale=0.8]{q1a.png} % Scales to 50% of original size
\end{center}
\noindent\rule{\textwidth}{0.4pt}\\
\newpage

%------------------
% Solution for 1(b)
%------------------
\subsection*{Solution 1}
\noindent\rule{\textwidth}{0.4pt}\\
\subsubsection*{Solution 1 (b)}
\subsubsection*{Python Code}
\begin{lstlisting}
import torch
from torchvision.utils import save_image
import matplotlib.pyplot as plt

# Set latent dim and number of steps
latent_dim = 4
num_steps = 256  
model.eval()     

with torch.no_grad():
    # Change range(20) to range(2) to get exactly 2 images
    for idx in range(2):
        # Pick random start and end points
        start_point = torch.randn(1, latent_dim).to(device)
        end_point = torch.randn(1, latent_dim).to(device)

        # Interpolate 256 points between start and end
        # This creates a smooth transition vector
        alphas = torch.linspace(0, 1, num_steps).to(device).unsqueeze(1).unsqueeze(2)
        interpolated_points = start_point + alphas * (end_point - start_point)
        # Note: shape is [256, 1, 4]

        # Decode these points into images
        generated_image = model.decoder(interpolated_points)
        
        # Save the grid
        # nrow=16 makes it a square 16x16 grid
        filename = f'./samples/latent_exploration_z_{latent_dim}_{idx+1}.png'
        save_image(generated_image.view(num_steps, 1, 28, 28), 
                   filename, nrow=16)
        print(f"Saved {filename}")

# Display the first one to check
img = plt.imread(f'./samples/latent_exploration_z_{latent_dim}_1.png')
plt.title(f'Latent space exploration (latent dim {latent_dim})')
plt.axis('off')
plt.imshow(img)
plt.show()
\end{lstlisting}

\subsubsection*{Generated Image}
\begin{center}
          
\includegraphics[scale=0.5]{./samples/latent_exploration_z_4_1.png} % Scales to 50% of original size
\end{center}
\begin{center}
          
\includegraphics[scale=0.5]{./samples/latent_exploration_z_4_2.png} % Scales to 50% of original size
\end{center}

\noindent\rule{\textwidth}{0.4pt}\\
\newpage

%------------------
% Solution for 2(a)
%------------------
\subsection*{Solution 2}
\noindent\rule{\textwidth}{0.4pt}\\
\subsubsection*{Solution 2 (a)}

To construct the 100-dimensional embeddings, the Positive Pointwise Mutual Information factorization approach was applied to the Brown corpus from the NLTK library. The steps are outlined below:

\begin{enumerate}
    \item \textbf{Vocabulary ($V$) and Context ($C$):} The top 5,000 most frequent words from the corpus to form the vocabulary $V$. For the context words $C$, the top 1,000 most frequent words were selected
    \item \textbf{Window Definition:} The context window was defined as a symmetric window of size 2 (i.e., 2 words preceding and 2 words succeeding the target word).
    \item \textbf{Co-occurrence Matrix Construction:} A $5000 \times 1000$ raw count matrix $M$ where $M_{ij}$ represents the number of times context word $c_j$ appeared within the window of target word $w_i$.
    \item \textbf{PPMI Reweighting:} $M$ was converted to a PPMI matrix where each entry is calculated as $\max(\log(\frac{P(w,c)}{P(w)P(c)}), 0)$. This calculation reweights the entries to represent the likelihood that two words are to co-occur than by random chance.
    \item \textbf{Dimensionality Reduction (SVD):} Singular Value Decomposition (SVD) on the sparse PPMI matrix to obtain the rank-100 approximation: $M_{PPMI} \approx U \Sigma V^T$. 
\end{enumerate}

\subsubsection*{\normalfont}{$\therefore$ The $\mathbb{R}^{100}$ is just the representation of the dimensionality reduction of the original $\mathbb{R}^{5000}$ space.}
\noindent\rule{\textwidth}{0.4pt}\\
\newpage

%------------------
% Solution for 2(b)
%------------------
\subsection*{Solution 2}
\noindent\rule{\textwidth}{0.4pt}\\
\subsubsection*{Solution 2 (b)}

Using Cosine Distance ($1 - \text{similarity}$), the nearest neighbor was found for 20 test words.  Below is a subset of the results:

\begin{itemize}
    \item \textbf{School} $\to$ \textit{College} (Related institution)
    \item \textbf{Music} $\to$ \textit{Jazz} (Category/Instance relationship)
    \item \textbf{Government} $\to$ \textit{Administration} (Synonym)
    \item \textbf{War} $\to$ \textit{Military} (Thematic relation)
    \item \textbf{He} $\to$ \textit{She} (Syntactic role/pronoun)
    \item \textbf{Book} $\to$ \textit{Novel} (Synonym)
    \item \textbf{Car} $\to$ \textit{Driver} (Functional relationship)
    \item \textbf{Water} $\to$ \textit{Air} (Elemental relationship)
    \item \textbf{Night} $\to$ \textit{Day} (Antonym/Temporal pair)
    \item \textbf{Love} $\to$ \textit{Life} (Thematic association)
\end{itemize}

\subsection*{\normalfont}{$\therefore$ The embedding successfully captures semantic proximity.}

\noindent\rule{\textwidth}{0.4pt}\\
\newpage

%------------------
% Solution for 2(c)
%------------------
\subsection*{Solution 2}
\noindent\rule{\textwidth}{0.4pt}\\
\subsubsection*{Solution 2 (c)}

The K-Means Clustering algorithm to group the 100-dimensional embeddings into $k=100$ clusters. 

\subsubsection*{Algorithm and Distance Choice}
\begin{itemize}
    \item \textbf{Algorithm: K-Means.} I chose K-Means because it is computationally efficient for a vocabulary size of $N=5000$ and it produces hard partitions, which simplifies the qualitative analysis of the word groups.
    \item \textbf{Distance Function: Euclidean (on L2-normalized vectors).} Standard K-Means optimizes for Euclidean distance. Euclidean distance and Cosine distance are monotonically related by:
    $$||x-y||^2 = 2(1 - \cos\theta)$$
\end{itemize}

\subsubsection*{Cluster Analysis:}
The resulting clusters exhibit high coherence, below are examples of this:

\begin{itemize}
    \item \textbf{Cluster 1 (Time/Duration):} \textit{day, year, week, month, hour, moment, night, morning}
    \item \textbf{Cluster 2 (Family/People):} \textit{mother, father, son, daughter, wife, husband, girl, boy}
    \item \textbf{Cluster 3 (Prepositions/Grammar):} \textit{in, on, at, from, by, with, for}
    \item \textbf{Cluster 4 (Politics/Civics):} \textit{government, state, nation, country, federal, law, policy}
\end{itemize}

\subsection*{\normalfont}{$\therefore$ The clusters are meaningful and distinct.}
\noindent\rule{\textwidth}{0.4pt}\\


\end{document}
